 \documentclass[11pt]{beamer}
\usepackage{amsfonts,amsmath,amsthm,amssymb}
\theoremstyle{plain}
\newtheorem{conjecture}{Conjecture}[section]
\usepackage{mathtools,mathptmx,listings,forest,enumitem}
\usepackage{graphicx}
\usepackage{pgfplots}
\pgfplotsset{compat=newest}
% plotting things
\usepackage{graphicx}
\graphicspath{{images/}}
\usepackage{tikz-cd}
\pgfplotsset{compat=1.15}
\usepackage[
	backend=biber,
	style=verbose,
	sorting=ynt
]{biblatex}
\addbibresource{references.bib}
\usetheme{Madrid}
\usepackage{float,mathtools,dirtytalk,ulem,csquotes,cancel,hyperref}
\usepackage{forest}
\usepackage{tikz-qtree}

\usepackage{tcolorbox}
\usepackage{subcaption}
\usepackage{quiver}

\author[] % (optional)
{Emon Hossain\inst{1}}

\institute[University of Dhaka] % (optional)
{
  \inst{1}%
  Lecturer\\MNS department\\Brac University
}

\date[] % (optional)
{\textsc{Lecture-09}}


\title[]{MAT215: Complex Variables And Laplace Transformations}

\setbeamertemplate{navigation symbols}{}


\AtBeginSection[]
{
  \begin{frame}
    \frametitle{Table of Contents}
    \tableofcontents[currentsection]
  \end{frame}
}

\usepackage{Kyushu}

% \usetheme{Frankfurt}

\begin{document}

\begin{frame}
\titlepage
\end{frame}

\begin{frame}{Definition of Harmonic Functions}
\begin{definition}
A function $u(x,y)$ is called a \textit{harmonic function} if it satisfies the Laplace's equation:
\[\nabla^2 u=\frac{\partial^2 u}{\partial x^2} + \frac{\partial^2 u}{\partial y^2} = 0.\]
\end{definition}
\end{frame}

\begin{frame}{Example}
\begin{example}
Show that the function $u(x,y) = 3x^2y+2x^2-y^3-2y^2$ is a harmonic function.
\end{example}
\begin{example}
Show that the function $u(x,y) = \tan^{-1}\left(\frac yx\right)$ is a harmonic function.
\end{example}
\end{frame}
\begin{frame}{Theorem}
    \begin{theorem}
    If $f(z) = u(x,y) + iv(x,y)$ is an analytic function in a region $D$, then both $u$ and $v$ are harmonic functions if they have continuous second partial derivatives in $D$.
    \end{theorem}
\end{frame}

\begin{frame}{Example}
\begin{example}
Show that the function $u(x,y) = 3x^2y+2x^2-y^3-2y^2$ is a harmonic function and find its harmonic conjugate $v(x,y)$.
Such that $f(z) = u(x,y) + iv(x,y)$ is an analytic function.
\end{example}
\end{frame}

\begin{frame}{Example}
\begin{example}
Show that the function $u(x,y) = e^{-2x}\sin(2y)$ is a harmonic function and find its harmonic conjugate $v(x,y)$.
Such that $f(z) = u(x,y) + iv(x,y)$ is an analytic function.
\end{example}
\end{frame}

\begin{frame}{Example}
\begin{example}
Show that the function $u(x,y) = \ln(x^2+y^2)$ is a harmonic function and find its harmonic conjugate $v(x,y)$.
Such that $f(z) = u(x,y) + iv(x,y)$ is an analytic function.
\end{example}
\end{frame}

\begin{frame}{Example}
\begin{example}
Show that the function $u(x,y) = \cos(x)\sinh(y)$ is a harmonic function and find its harmonic conjugate $v(x,y)$.
Such that $f(z) = u(x,y) + iv(x,y)$ is an analytic function.
\end{example}
\textbf{Hint: }
$$\frac{d}{dx}(\sinh x)=\cosh x, \quad \frac{d}{dx}(\cosh x)=\sinh x$$
\end{frame}

\begin{frame}{Example}
\begin{example}
Show that the function $u(x,y) = xe^{-x}\sin(y)-e^{-x}y\cos(y)$ is a harmonic function and find its harmonic conjugate $v(x,y)$.
Such that $f(z) = u(x,y) + iv(x,y)$ is an analytic function.
\end{example}
\end{frame}
\end{document}